{\scriptsize
\begin{minipage}{0.5\textwidth}
    \begin{tabularx}{\linewidth}{|X|X|}
        \hline \multicolumn{1}{|c|}{ Notions et contenus } & \multicolumn{1}{c|}{ Capacités exigibles } \\
        \hline 1.1. Modèle scalaire des ondes lumineuses \\
        \hline 
        Modèle de propagation dans l'approximation de l'optique géométrique. & 
        Associer la grandeur scalaire de l'optique à une composante d'un champ électrique. \\
        \hline 
        Vibration lumineuse. Chemin optique. Déphasage dû à la propagation. & 
        Exprimer le retard de phase en un point en fonction de la durée de propagation ou du chemin optique. \\
        \hline 
        Surfaces d'ondes. Théorème de Malus. Onde plane, onde sphérique ; effet d'une lentille mince dans l'approximation de Gauss. & 
        Utiliser l'égalité des chemins optiques sur les rayons d'un point objet à son image. Associer une description de la formation des images en termes de rayons lumineux et en termes de surfaces d'onde. \\
        \hline
    \end{tabularx}
\end{minipage}%
\begin{minipage}{0.5\textwidth}
    \begin{tabularx}{\linewidth}{|p{3cm}|X|}
        \hline Modèle d'émission. & \\
        \hline Réception d'une onde lumineuse. & 
        Largeur spectrale. Cohérence temporelle. 
        Classer différentes sources lumineuses 
        (lampe spectrale basse pression, laser, 
        source de lumière blanche..) en fonction du 
        temps de cohérence de leurs diverses 
        radiations. 
        Citer quelques ordres de grandeur des 
        longueurs de cohérence temporelle 
        associées à différentes sources. 
        Relier, en ordre de grandeur, le temps de 
        cohérence et la largeur spectrale de la 
        radiation considérée. \\
        \hline Récepteurs. Intensité lumineuse. & 
        Comparer le temps de réponse d'un 
        récepteur usuel (œil, photodiode, capteur 
        CCD) aux temps caractéristiques des 
        vibrations lumineuses. 
        Relier l'intensité lumineuse à la moyenne 
        temporelle du carré de la grandeur scalaire 
        de l'optique. \\
        \hline
    \end{tabularx}
\end{minipage}
}
